\section{Cambio del cuerpo de base en los preesquemas algebraicos}
\label{IV.2.4-fr}

\subsection{Dimensión de los preesquemas algebraicos.}
\label{subsection:IV.2.4.1-fr}

Desarrollaremos en el \S~5 la teoría general de la dimensión de los
preesquemas; pero la teoría de la dimensión de los preesquemas
\emph{algebraicos} puede desarrollarse de forma más elemental y, como presenta
además numerosos caracteres especiales, daremos aquí una exposición rápida,
independiente de la teoría general.

Si $K$ es un cuerpo y $L$ una extensión de $K$, se denotará por
$\operatorname{deg.tr}_K L$ el \emph{grado de trascendencia} de $L$ sobre $K$.

\begin{definition}[4.1.1]
\label{IV.2.4.1.1-fr}
Sean $k$ un cuerpo, $X$ un preesquema localmente de tipo finito sobre $k$. Se
llama dimensión de $X$ al número
\begin{equation}
\tag{4.1.1.1}
\label{IV.2.4.1.1.1-fr}
\dim X=\sup_x\operatorname{deg.tr}_k \boldsymbol{k}(x)
\end{equation}
donde $x$ recorre el conjunto de los puntos maximales de $X$.
\end{definition}

Veremos en el \S~5 (5.2.2) que la definición (4.1.1) solo depende en
apariencia del cuerpo de base $k$ sobre el cual $X$ es localmente de tipo
finito, y que el número $\dim(X)$ así definido coincide con la dimensión del
\emph{espacio subyacente} $X$ (definida en $(\mathbf{0},14.1.2)$). Obviamente se tiene
$\dim X=\dim X_{\mathrm{red}}$.

Se observará que cada $\boldsymbol{k}(x)$ es una extensión \emph{de tipo finito} de $k$
(\textbf I, 6.3.3), por tanto tiene grado de trascendencia finito sobre $k$. Si $X$ es
de tipo finito sobre $k$ y no vacío, es noetheriano (\textbf I, 6.3.7), por lo que solo
tiene un número finito de componentes irreducibles, y en consecuencia
$\dim(X)$ es \emph{finita} y $\geqslant0$; la definición (4.1.1) da
\[
\dim(\varnothing)=\setmin\infty.
\]

Es claro que, si se denota por $(X_\alpha)$ la familia de los subpreesquemas
cerrados reducidos de $X$ cuyos espacios subyacentes son las componentes
irreducibles de $X$ (\textbf I, 5.2.1), se tiene
\begin{equation}
\tag{4.1.1.2}
\label{IV.2.4.1.1.2-fr}
\dim(X)=\sup_\alpha\dim(X_\alpha).
\end{equation}
Esto reduce, pues, el cálculo de la dimensión al caso de los preesquemas
\emph{íntegros} (localmente de tipo finito sobre $k$). Finalmente, se tiene
obviamente
\begin{equation}
\tag{4.1.1.3}
\label{IV.2.4.1.1.3-fr}
\dim(X)=\dim(U)
\end{equation}
para todo subpreesquema inducido sobre un abierto \emph{en todas partes denso}
$U$ de $X$; esto reduce, pues, en último término la noción de dimensión al
caso de un esquema \emph{afín} de tipo finito sobre $k$.

\begin{theorem}[4.1.2]
\label{IV.2.4.1.2-fr}
Sean $X$, $Y$ dos preesquemas localmente de tipo finito sobre un cuerpo $k$,
$f:X\to Y$ un $k$-morfismo.

\oldpage[IV-2]{53}

\begin{enumerate}
\item[\textup{(i)}] Si $f$ es cuasicompacto y dominante, se tiene:
$\dim(Y)\leqslant\dim(X)$.
\item[\textup{(ii)}] Si $f$ es cuasifinito, se tiene:
$\dim(X)\leqslant\dim(Y)$.
\item[\textup{(iii)}] Supongamos que $X$ es de tipo finito sobre $k$. Para que
$\dim(X)\geqslant n$ \textup{(resp. $\dim(X)\leqslant n$, $\dim(X)=n$)}, es necesario y
suficiente que exista un abierto $U$ denso en $X$, y un $k$-morfismo
$g:U\longrightarrow\mathbf V(k^n)$
$(=\operatorname{Spec}(k[T_1,\ldots,T_n])$,
que denotaremos también por $\mathbf V_k^n)$ que sea sobreyectivo \textup{(resp.} finito,
\textup{resp.} finito y sobreyectivo\textup{)}. Si $X$ es un esquema afín, se puede tomar $U=X$.
\end{enumerate}
\end{theorem}

\begin{proof}
(i) Sea $y$ un punto maximal de $Y$; se sabe que $f^{-1}(y)$ contiene un
punto maximal $x$ de $X$ (1.1.5), luego $\boldsymbol{k}(x)$ es una extensión de $\boldsymbol{k}(y)$,
de donde se obtiene la desigualdad
$\operatorname{deg.tr}_k \boldsymbol{k}(y)
\leqslant\operatorname{deg.tr}_k \boldsymbol{k}(x)
\leqslant\dim(X)$,
lo que demuestra (i).

(ii) Si $x$ es un punto maximal de $X$, $\boldsymbol{k}(x)$ es una extensión finita de
$\boldsymbol{k}(f(x))$ (\textbf{II}, 6.2.2), por tanto tiene el mismo grado de trascendencia sobre
$k$. Si se considera el subpreesquema cerrado reducido de $Y$ cuyo espacio
subyacente es $\overline{\{f(x)\}}$, se ve que basta demostrar el
\end{proof}

\begin{corollary}[4.1.2.1]
\label{IV.2.4.1.2.1-fr}
Sea $Y$ un $k$-preesquema localmente de tipo finito; para todo subpreesquema
$Z$ de $Y$, se tiene $\dim(Z)\leqslant\dim(Y)$. Supongamos además que las
componentes irreducibles de $Y$ tienen todas la misma dimensión; entonces, para
que $Z$ sea raro en $Y$, es necesario y suficiente que
$\dim(Z)<\dim(Y)$.
\end{corollary}

\begin{proof}
En efecto, sea $z$ un punto maximal de $Z$; considerando uno de los
subpreesquemas cerrados reducidos de $Y$ cuyo espacio subyacente sea una de las
componentes irreducibles de $Y$ que contiene a $z$, la cuestión se reduce al caso en
que $Y$ es íntegro con punto genérico $y$, y luego, considerando en $Y$ un
abierto afín que contiene a $z$, al caso en que $Y$ es afín y de tipo finito
sobre $k$, es decir, $Y=\operatorname{Spec}(A)$, y $Z$ es cerrado en $Y$, esto
es, $Z=\operatorname{Spec}(A/\mathfrak a)$, donde $\mathfrak a$ es un ideal
del anillo íntegro $A$, distinto de $A$. En virtud del lema de normalización
(Bourbaki, \emph{Alg. comm.}, cap.~V, \S~3,
n\textsuperscript{o}~1, th.~1), existe una sucesión finita
$(x_i)_{1\leqslant i\leqslant n}$ de elementos de $A$, algebraicamente
independientes sobre $k$, tales que $A$ es entero sobre el anillo
$B=k[x_1,\ldots,x_n]$ y que $\mathfrak a\cap B$ está generado por una
subfamilia $(x_i)_{1\leqslant i\leqslant p}$ (posiblemente vacía) de
$(x_i)_{1\leqslant i\leqslant n}$. Sea
$g:Y\to\operatorname{Spec}(B)=\mathbf V_k^n$ el morfismo finito dominante que
corresponde a la inyección canónica $B\to A$; como
$C=B/(\mathfrak a\cap B)$ es isomorfo a $k[x_{p+1},\ldots,x_n]$, $g$ induce
en $Z$ un morfismo finito dominante $Z\to\mathbf V_k^{n-p}$; se tiene, pues,
$\operatorname{deg.tr}_k \boldsymbol{k}(y)=n$ y
$\operatorname{deg.tr}_k \boldsymbol{k}(z)=n\setmin p$, ya que $\boldsymbol{k}(y)$ (resp. $\boldsymbol{k}(z)$) es una
extensión finita de $\boldsymbol{k}(g(y))$ (resp. $\boldsymbol{k}(g(z))$). Esto demuestra la primera
afirmación de (4.1.2.1). Además, si $p=0$, necesariamente $z=y$, pues $y$ es
el único punto de $Y$ cuya imagen por $g$ es el punto genérico de
$\mathbf V_k^n$; en este caso, $Z$ contiene, por tanto, un abierto no vacío de
$Y$. Si, por el contrario, $p>0$, necesariamente $z\neq y$ y
$\overline{\{z\}}$ es, por tanto, \emph{raro} en $\overline{\{y\}}$, lo que concluye
la demostración de (4.1.2.1).
\end{proof}

\begin{proof}
(iii) El hecho de que las condiciones enunciadas sean suficientes resulta
inmediatamente de (i) y (ii) y de (1.5.4, (v)). Para demostrar que son
necesarias, se puede considerar un abierto $U$ denso en $X$, unión de abiertos
afines $U_i$ dos a dos disjuntos e irreducibles $(1\leqslant i\leqslant m)$,
cada uno de los cuales contiene uno de los puntos maximales de $X$
$(\mathbf{0}_{\mathrm I},2.1.6)$; además, se puede suponer que $X$ es reducido y, si
$X$ es afín, se puede naturalmente tomar $U=X$. El mismo

\oldpage[IV-2]{54}

razonamiento que en (4.1.2.1) muestra que si $\dim(U_i)=n_i$, existe un
$k$-morfismo
$g_i:U_i\to\mathbf V_k^{n_i}$ que es finito y dominante, por tanto
sobreyectivo (\textbf{II}, 6.1.10). Sea $n=\dim(X)=\sup_i n_i$; para cada $n_i$ hay,
por tanto, un morfismo $h_i:\mathbf V_k^{n_i}\to\mathbf V_k^n$ que es una
inmersión cerrada correspondiente al homomorfismo canónico
$k[T_1,\ldots,T_n]\to k[T_1,\ldots,T_{n_i}]$, y que es la identidad para
$n_i=n$. Como $U$ es suma de los $U_i$, se tomará para $g$ el morfismo que en
cada $U_i$ coincide con $h_i\circ g_i$; es evidentemente finito y
sobreyectivo. Cuando $n'\geqslant n$, se tendrá un morfismo finito
$U\to\mathbf V_k^{n'}$ componiendo la inmersión cerrada canónica
$\mathbf V_k^n\to\mathbf V_k^{n'}$ con $g$; cuando $n'\leqslant n$, se compone
de igual modo con $g$ el morfismo canónico
$p:\mathbf V_k^n\to\mathbf V_k^{n'}$ correspondiente a la inyección canónica
$k[T_1,\ldots,T_{n'}]\to k[T_1,\ldots,T_n]$, observando que $p$ es fielmente
plano, por tanto \emph{sobreyectivo}.
\end{proof}

\begin{remark}[4.1.3]
\label{IV.2.4.1.3-fr}
El corolario (4.1.2.1) muestra que en la fórmula (4.1.1.1), se puede suponer
que $x$ recorre el conjunto de \emph{todos los puntos} de $X$.
\end{remark}
